\chapter*{Resumo}

A gestão financeira pessoal é um desafio crescente na sociedade moderna, agravado pela multiplicidade de contas bancárias, produtos financeiros e a complexidade do sistema fiscal português. As soluções existentes no mercado, maioritariamente de origem norte-americana, não atendem às especificidades do mercado português, nomeadamente a integração com o sistema fiscal (IRS), a rede SIBS/Multibanco e os hábitos de consumo nacionais.

Este projeto propõe o desenvolvimento de uma plataforma web de gestão financeira pessoal especificamente concebida para o mercado português. A plataforma integra Open Banking (PSD2) através da API Salt Edge para agregação automática de contas bancárias, categorização inteligente de transações com recurso a Machine Learning (scikit-learn) e um simulador de IRS baseado nas regras fiscais portuguesas vigentes.

A arquitetura do sistema segue um modelo de microserviços, utilizando React 18 com TypeScript no frontend, Node.js com Express no backend e PostgreSQL como base de dados principal. O design da interface segue o paradigma Liquid Glass, priorizando a usabilidade e acessibilidade (WCAG AA).

A metodologia de desenvolvimento adotada é Scrum adaptada para desenvolvimento individual, organizada em 12 sprints de duas semanas, abrangendo o período de abril a setembro de 2026. A avaliação do sistema inclui testes de usabilidade com 10 a 15 utilizadores reais, utilizando o questionário SUS (System Usability Scale), complementados por métricas de performance técnica e precisão do modelo de Machine Learning.

\vspace{0.5cm}
\noindent\textbf{Palavras-chave:} Gestão Financeira Pessoal, Open Banking, PSD2, Machine Learning, IRS, Mercado Português, React, Node.js, Liquid Glass

\chapter*{Abstract}

Personal financial management is a growing challenge in modern society, exacerbated by multiple bank accounts, financial products, and the complexity of the Portuguese tax system. Existing market solutions, mostly of North American origin, fail to address the specificities of the Portuguese market, namely integration with the tax system (IRS), the SIBS/Multibanco network, and national consumption habits.

This project proposes the development of a personal financial management web platform specifically designed for the Portuguese market. The platform integrates Open Banking (PSD2) through the Salt Edge API for automatic bank account aggregation, intelligent transaction categorization using Machine Learning (scikit-learn), and an IRS simulator based on current Portuguese tax rules.

The system architecture follows a microservices model, using React 18 with TypeScript on the frontend, Node.js with Express on the backend, and PostgreSQL as the main database. The interface design follows the Liquid Glass paradigm, prioritizing usability and accessibility (WCAG AA).

The development methodology adopted is Scrum adapted for individual development, organized in 12 two-week sprints, covering the period from April to September 2026. System evaluation includes usability tests with 10 to 15 real users, using the SUS (System Usability Scale) questionnaire, complemented by technical performance metrics and Machine Learning model accuracy.

\vspace{0.5cm}
\noindent\textbf{Keywords:} Personal Financial Management, Open Banking, PSD2, Machine Learning, IRS, Portuguese Market, React, Node.js, Liquid Glass
